% Exact proof module, 6 September 2026. Input after es_niemeier_shared.tex.
% All original coordinates and Euclidean root lengths are retained.
\section{The exact Leech neighbor of the retained index-72 lattice}
\label{esnl:sec:neighbor}

Retain the lattice, metric, discriminant marking, and glue literally:
\[
 \begin{split}
 A&=\{(x_0,\ldots,x_5)\in\mathbb Z^6:\sum x_j=0\},\\
 D&=\{(d_1,\ldots,d_4)\in\mathbb Z^4:\sum d_i\equiv0\pmod2\},
 \qquad R=A^{\perp4}\perp D,\\
 C_2&=\langle b_0=(1111;00),b_1=(1100;10),b_2=(1010;01)\rangle_{\mathbb F_2},\\
 C_3&=\langle z_1=(1,1,1,0),z_2=(1,2,0,1)\rangle_{\mathbb F_3},\\
 H&=\{(3u+2z;y):(u;y)\in C_2,\ z\in C_3\},\\
 N&=\{x\in R^\vee:x+R\in H\}.
 \end{split}
\]
The coordinates in each $A$ remain the six ordered elements $f_i289^j$
of its original residue fibre. The $A$ discriminant generator is
$\lambda=(5,-1,-1,-1,-1,-1)/6$. The $D$ markings $10,01,11$
remain represented by $v_D=(1,0,0,0)$,
$s_D=(1,1,1,1)/2$, $c_D=(-1,1,1,1)/2$.
The preceding proof establishes that $N$ is even unimodular,
$[N:R]=72$, $\det R=5184$, and its roots are exactly those of
$A_5^4D_4$. Every root has squared norm two. No metric rescaling
occurs anywhere below.

\subsection{The Weyl vector and the exact neighbor}

Use the original simple roots $a_j=e_j-e_{j+1}$, $0\le j\le4$,
in each $A$ component and
\[
 d_1=e_1-e_2,\quad d_2=e_2-e_3,\quad
 d_3=e_3-e_4,\quad d_4=e_3+e_4
\]
in $D$. Their highest roots are
$\theta_A=e_0-e_5=\sum_ja_j$ and
$\theta_D=e_1+e_2=d_1+2d_2+d_3+d_4$.
Both highest-root heights are five, so the common Coxeter number
is six. Define
\[
 \rho_A=(5,3,1,-1,-3,-5)/2,\qquad
 \rho_D=(3,2,1,0),\qquad \rho=\rho_A^{\oplus4}\perp\rho_D.
 \label{esnl:eq:rho}
\]
Each simple root pairs to one with $\rho$.
Moreover $\rho_A-3\lambda=(0,2,1,0,-1,-2)\in A$ and
$\rho_D\in D$. Thus the exact glue class of $\rho$ is
$(3,3,3,3;00)$, the image of $b_0$, and $\rho\in N$.
Directly
\[
 \rho_A^2=35/2,\quad \rho_D^2=14,\quad
 \rho^2=4(35/2)+14=84=2\cdot6\cdot7.
\]
Consequently $(\rho,N)\subset\mathbb Z$.
Fix $\alpha=a_0$ in the first $A$ component and put
\[
 \ell_\rho:N\to\mathbb Z/6,\quad x\mapsto(\rho,x)\bmod6,\qquad
 K=\ker\ell_\rho,\quad R_0=R\cap K,\quad
 v=\rho/6-\alpha,\quad \Lambda=K+\mathbb Z v.
 \tag{NL1}
\]
The pairing $(\rho,\alpha)=1$ proves
$[N:K]=[R:R_0]=6$. For $x\in K$, the number
$(v,x)=(\rho,x)/6-(\alpha,x)$ is integral, and
\[
 v^2=84/36-2/6+2=4,\qquad
 6v=\rho-6\alpha\in K,\qquad(\rho,6v)=78.
\]
Hence $\Lambda$ is positive definite even integral of rank 24.
If $mv\in N$, its pairing after adding $m\alpha$ with $\alpha$
is $m/6\in\mathbb Z$, so $6\mid m$.
It follows that
\[
 N\cap\Lambda=K,\qquad [\Lambda:K]=6,\qquad
 \det K=36,\qquad \det\Lambda=36/6^2=1.
 \tag{NL2}
\]
Changing $\alpha$ to another root of $\rho$-pairing one changes
$v$ by an element of $K$, and so does not change $\Lambda$.
We shall use the equivalent exact congruence description
\[
 \Lambda=\{x+k\rho/6:x\in N,\ k\in\mathbb Z,\
                                (\rho,x)+k\equiv0\pmod6\}.
 \tag{NL3}
\]
Indeed write $x+k\rho/6=(x+k\alpha)+kv$; the first summand
is in $K$ exactly under this congruence. Changing $k$ by six
can be absorbed in $x$ using $\rho\in N$, and preserves the
congruence since $\rho^2=84$.

\subsection{A complete rootlessness proof}

\begin{theorem}[The Leech neighbor in the retained coordinates]
\label{esnl:thm:leech-neighbor}
The lattice $\Lambda=K+\mathbb Z(\rho/6-\alpha)$ is positive
definite even unimodular of rank 24 and has minimum squared norm
exactly four. It is the holy-construction Leech lattice of the
retained $A_5^4D_4$ glue, and
$N\cap\Lambda=K$ with $[N:K]=[\Lambda:K]=6$.
\end{theorem}

\begin{proof}
The rank, parity, unimodularity and indices were proved in
(NL1)--(NL2). The holy-presentation identification is proved in
(NL8)--(NL11) below. We now prove the required minimum directly.
For $k=0$ in (NL3), a possible root is an original root of $N$.
An $A$ root pairs with $\rho$ to one of $\pm1,\ldots,\pm5$.
For $D$, all roots $\pm e_i\pm e_j$ have nonzero $\rho_D$-pairing
of absolute value at most five. Thus none belongs to $K$.
Negation and the reduction of $k$ modulo six leave $k=1,2,3$.

Define the exact shifted class minima by
\[
 m_A(k,t)=\min_{x\in t\lambda+A}|x+k\rho_A/6|^2,\qquad
 m_D(k,y)=\min_{x+D=y}|x+k\rho_D/6|^2.
\]
Their values are
\[
 \begin{array}{c|rrrrrr}
 k\backslash t&0&1&2&3&4&5\\\hline
 1&35/72&35/72&35/72&35/72&35/72&35/72\\
 2&11/18&4/9&11/18&4/9&11/18&4/9\\
 3&3/8&17/24&17/24&3/8&17/24&17/24
 \end{array}
 \quad
 \begin{array}{c|rrrr}
 k\backslash y&00&10&01&11\\\hline
 1&7/18&7/18&7/18&7/18\\
 2&2/9&5/9&5/9&5/9\\
 3&1/2&1/2&1/2&1/2
 \end{array}.
 \tag{NL4}
\]
We prove the table and its needed equality information.
The shifted $A$ coordinates are precisely
\[
 \left(n_j-\frac t6+\frac{k(5-2j)}{12}\right)_{j=0}^5,
               \qquad n_j\in\mathbb Z,\quad\sum n_j=t.
 \tag{NL5}
\]
At $k=1$, coordinatewise nearest representatives are the six
numbers $\pm1/12,\pm3/12,\pm5/12$ in an order depending on $t$.
They sum to zero and their squares sum to $35/72$.
At $k=2$ and even $t$, the nearest coordinates consist of
two copies of each sign of $1/6$ and one copy of each sign of
$1/2$; the half-integer ties admit exactly the needed zero-sum
choice and the squared sum is $11/18$.
At odd $t$, there are two each of $0,1/3,-1/3$, giving $4/9$.
For $t=1,3,5$ their actual ordered shifted vectors are
\[
 (-1,1,0,-1,1,0)/3,\quad
 (1,0,-1,1,0,-1)/3,\quad
 (0,-1,1,0,-1,1)/3.
 \tag{NL6}
\]
Subtracting $\rho_A/3$ gives original squared norms
$17/6,3/2,17/6$, respectively.
At $k=3,t=0,3$, the nearest coordinates alternate $\pm1/4$,
giving $3/8$ with zero sum.
At each other $t$, the coordinatewise nearest representatives
have three absolute values $1/12$ and three $5/12$, square sum
$13/24$, and coordinate sum $1$ or $-1$.
Meeting the zero-sum constraint requires an integer change.
Every nonzero coordinate change increases the squared sum by
at least $(7/12)^2-(5/12)^2=1/6$.
Changing one suitable $5/12$ coordinate attains this increase
and restores the zero sum. Thus the constrained minimum is
$17/24$. At $t=0,3$, the original minimizing vectors are
\[
 (-1,-1,0,0,1,1),\qquad (-3,-1,-1,1,1,3)/2,
 \tag{NL7}
\]
of squared norms $4,11/2$.

For $D$ and $k=1$ the shift is $(1/2,1/3,1/6,0)$.
The nearest integer coordinate distances are
$(1/2,1/3,1/6,0)$; the first-coordinate tie meets either integer
parity. The nearest half-integer distances are
$(0,1/6,1/3,1/2)$; the last-coordinate tie meets either
half-integer class. Both squared sums are $7/18$.
At $k=2$ the shift is $(1,2/3,1/3,0)$.
The unique nearest integer vector before shifting is
$(-1,-1,0,0)\in D$, giving $2/9$.
Changing parity costs at least $1/3$, attained by changing the
second or third coordinate, giving $5/9$.
The half-integer coordinate distances are
$(1/2,1/6,1/6,1/2)$, giving $5/9$; ties meet either marking.
At $k=3$ the shift is $(3/2,1,1/2,0)$.
Integer-class minimum distances are $(1/2,0,1/2,0)$,
and half-integer-class distances $(0,1/2,0,1/2)$.
Ties meet every marking and all square sums are $1/2$.
In class $00$, the two original minimizing vectors are
\[
 (-1,-1,0,0),\qquad(-2,-1,-1,0),
\]
of squared norms $2,6$. In class $10$, they are
\[
 (-1,-1,-1,0),\qquad(-2,-1,0,0),
\]
of squared norms $3,5$.
In either half-integer class they have the shapes
$(-3/2,-1/2,-1/2,\pm1/2)$ and
$(-3/2,-3/2,-1/2,\pm1/2)$.
Exactly one last-coordinate sign in each shape meets each
original marking; their squared norms are $3,5$.
This proves every entry and the relevant equality cases in (NL4).

For $k=1$, every full shifted vector has squared norm at least
$4(35/72)+7/18=7/3>2$.
For $k=2$, the lower bound from (NL4) is two.
Equality forces all $t_i$ odd and $y=00$.
In the actual glue this forces $u=1111$.
If $z=0$, all $t_i=3$ and (NL6) gives
$x^2=4(3/2)+2=8$.
If $z\ne0$, its weight is three, so exactly one $t_i=3$
and three $t_i\in\{1,5\}$; then
$x^2=3/2+3(17/6)+2=12$.
The norm-two equation is
\[
 2=x^2+\tfrac23(\rho,x)+\tfrac{28}{3},
 \qquad(\rho,x)=-11-\tfrac32x^2.
\]
Its pairing is $-23$ or $-29$, always $1$ modulo six.
Condition (NL3) requires residue $4$, so all nine equality
candidates, one per ternary word, are excluded.

For $k=3$, the lower bound two requires $z=0$:
every nonzero ternary word raises the minimum by
$3(17/24-3/8)=1$.
If $u=0$, the possible original squared norms are $18,22$.
If $u=1111$, they are $24,28$.
If $u$ has weight two, its $D$ class is nonzero and they are
$2(4)+2(11/2)+3=22$ or $2(4)+2(11/2)+5=24$.
These are all sixteen equality candidates.
But
\[
 2=x^2+(\rho,x)+21,\qquad(\rho,x)=-19-x^2,
\]
and (NL3) would require $x^2\equiv2\pmod6$.
None of $18,22,24,28$ has that residue.
Thus $\Lambda$ has no roots. Together with (NL2), this proves
that it is positive definite even unimodular of rank 24,
with minimum squared norm exactly four, attained by $v$.
\end{proof}

\subsection{The exact holy-construction presentation}

For $0\le t\le5$ set
\[
 \omega_t=(\underbrace{1,\ldots,1}_{t},
             \underbrace{0,\ldots,0}_{6-t})-\tfrac t6(1,\ldots,1).
\]
Then $\omega_t+A=t\lambda+A$ and
$(\rho_A,\omega_t)=t(6-t)/2=3\omega_t^2$.
For the $D$ markings set
\[
 \omega_{00}=0,\quad\omega_{10}=e_1,\quad
 \omega_{01}=(1,1,1,1)/2,\quad
 \omega_{11}=(1,1,1,-1)/2.
\]
The last differs from the original $c_D$ by $(1,0,0,-1)\in D$,
so retains its exact marking.
For $a=(t_1,t_2,t_3,t_4;y)\in H$ define
\[
 u_a=(\omega_{t_1},\omega_{t_2},\omega_{t_3},\omega_{t_4};\omega_y),
 \qquad g_a=u_a-\rho/6.
 \tag{NL8}
\]
These give all 72 representatives of $N/R$. Each nonzero $D$
representative has norm one and pairing three; hence
\[
 (\rho,u_a)=3u_a^2\in6\mathbb Z,\quad
 u_a\in K,\quad g_a^2=u_a^2-(\rho,u_a)/3+84/36=7/3.
 \tag{NL9}
\]
Let $\mathcal F$ be the 29 vectors consisting, in each original
component, of the negative original simple roots and the positive
highest root. Their pairings with $\rho$ are $-1$ or $5$.
There are positive integral zero relations of coefficient sum six:
\[
 \theta_A+\sum_{j=0}^4(-a_j)=0,\qquad
 \theta_D+(-d_1)+2(-d_2)+(-d_3)+(-d_4)=0.
 \tag{NL10}
\]
Every component of $g_a$ pairs to $1/6$ with the corresponding
vertices except one tip, where the pairing is $-5/6$.
For $A$ this follows because $\omega_t$ has simple-root pairing
one only at $a_{t-1}$ when $t\ne0$, and has
$(\omega_t,\theta_A)=1$ for $t\ne0$.
For $D$ the three nonzero representatives have simple-root
pairing one only at $d_1,d_4,d_3$, respectively; all have
highest-root pairing one. The zero representative points to
the highest root. These computations fix all orientations in
the holy generator diagram.

The following coefficient presentations are exact equalities:
\[
 \begin{split}
 N&=\left\{\sum_fa_ff+\sum_ab_ag_a:
                    a_f,b_a\in\mathbb Z,\ \sum_ab_a=0\right\},\\
 \Lambda&=\left\{\sum_fa_ff+\sum_ab_ag_a:
            a_f,b_a\in\mathbb Z,\ \sum_fa_f+\sum_ab_a=0\right\}.
 \end{split}
 \tag{NL11}
\]
For the first, the $f$ span $R$ and $g_a-g_0=u_a$ generate its
given extension $N$; the zero coefficient sum cancels $\rho/6$.
For the second, write $A=\sum a_f$, $B=\sum b_a$,
$r=\sum a_ff$, and $x=r+\sum b_au_a$.
The vector is $x-B\rho/6$, and by (NL9)
$(\rho,x)\equiv-A\equiv B$, exactly (NL3) with $k=-B$.
Conversely, for any $x+k\rho/6$ in (NL3), write
$x=r+\sum b_au_a$. The coefficient of $u_0=0$ can be set so
$B=-k$. Write $r$ with the negative simple-root vertices.
The congruence gives $A+B\equiv0\pmod6$.
Adding a suitable multiple of a zero relation (NL10) makes
$A+B=0$ exactly. Thus both directions of (NL11) are proved.

These are exactly the holy coefficient rules of Conway--Sloane
\cite[pp.~277--280]{esn:ConwaySloane1982}.
Borcherds gives the uniform proof of the holy-construction
theorem \cite[Section~7, especially the coefficient rules
(i)--(ii) and Lemma~7.3]{esn:Borcherds1985}.
Thus (NL11) is an explicit realization of that Leech
construction from the actual 72 marked glue words; all its
metric, determinant, integrality and rootlessness assertions
have also been independently proved here.

\subsection{The indices, discriminant morphism, and cyclic order twelve}

Reduction modulo $R$ gives an isomorphism $K/R_0\cong H$:
its kernel is $R_0$, and (NL9) supplies representatives in $K$
for every $H$ class. Consequently
\[
 [K:R_0]=72,\quad [N:R_0]=[\Lambda:R_0]=432,\quad
 \det R_0=6^2\cdot5184=186624.
 \tag{NL12}
\]
Thus the original and new inclusions, with every index, are
\[
 R_0\subset R\subset N,\quad [R:R_0]=6,\ [N:R]=72;
 \qquad
 R_0\subset K\subset\Lambda,\quad[K:R_0]=72,\ [\Lambda:K]=6.
\]
In particular $R\cap\Lambda=R_0$, an exact common sublattice.

\begin{theorem}[The exact order-twelve quotient generator]
\label{esnl:thm:cyclic-twelve}
The quotient of $\Lambda$ by the original height-zero root
sublattice $R_0$ is
$\Lambda/R_0\cong\mathbb Z/12\oplus(\mathbb Z/6)^2$.
The class of $v=\rho/6-\alpha$ has exact order twelve.
\end{theorem}

\begin{proof}
Let $\bar b_i$ denote the image in $H$ of the binary generator
$b_i$, namely its first four coordinates multiplied by three
and its last two coordinates retained.
Since $6v=\rho-6\alpha$ and $6\alpha\in R_0$,
\[
 6[v]=\bar b_0\quad\text{in }\Lambda/R_0.
 \tag{NL13}
\]
This is a nonzero order-two element. The classes
$\bar b_0,\bar b_1,\bar b_2,2z_1,2z_2$ form the displayed
independent binary and ternary generators of $H$.
Eliminate $\bar b_0$ using (NL13). The resulting surjective
presentation is
\[
 \Lambda/R_0\cong
 \mathbb Z/12\oplus(\mathbb Z/2)^2\oplus(\mathbb Z/3)^2
 \cong\mathbb Z/12\oplus(\mathbb Z/6)^2.
 \tag{NL14}
\]
There are no further relations: this presentation has order
$12\cdot4\cdot9=432$, equal to (NL12).
Thus $[v]$ has exact order twelve. This particular occurrence
of twelve is proved from the Coxeter denominator six and the
actual order-two Weyl-vector glue class.
\end{proof}

Set $M=N+\Lambda=N+\mathbb Z(\rho/6)$.
Then
\[
 M=K^\vee,\qquad K^\vee/K=
    \langle[\alpha],[v]\rangle\cong(\mathbb Z/6)^2.
 \tag{NL15}
\]
Indeed $M\subset K^\vee$, its index over $N$ is six by pairing
with $\alpha$, and $[K^\vee:N]=[N:K]=6$.
The cyclic images of $N/K$ and $\Lambda/K$ have order six and
trivial intersection, proving the stated generators.
Their exact discriminant data are
\[
 q_K([\alpha])=q_K([v])=0,\quad
 b_K([\alpha],[v])=1/6\pmod{\mathbb Z},\quad
 q_K(a[\alpha]+b[v])=ab/3\pmod{2\mathbb Z}.
 \tag{NL16}
\]
These follow from $\alpha^2=2$, $v^2=4$ and
$(\alpha,v)=1/6-2=-11/6$.
The two original lattices are therefore the inverse images
of the two indicated order-six isotropic subgroups in this
specified discriminant form.

\subsection{An exact Lorentzian quotient and symmetry transport}

Take $\mathcal I=N\perp U$ in coordinates $(x,m,n)$ with
\[
 \langle(x,m,n),(x',m',n')\rangle
 =(x,x')-mn'-m'n .
\]
Thus $U$ has Gram matrix
$\left(\begin{smallmatrix}0&-1\\-1&0\end{smallmatrix}\right)$.
The vectors $z=(0,0,1)$ and $w=(\rho,6,7)$ are primitive
isotropic: $w^2=84-2\cdot6\cdot7=0$, and six and seven are
coprime. There are integral quotient isometries
\[
 \begin{split}
 z^\perp/\mathbb Zz&\longrightarrow N,\quad[(x,0,n)]\mapsto x,\\
 w^\perp/\mathbb Zw&\longrightarrow\Lambda,\quad
                                  [(x,m,n)]\mapsto x-m\rho/6.
 \end{split}
 \tag{NL17}
\]
For the second map, orthogonality means $(\rho,x)=7m+6n$,
giving (NL3) with $k=-m$. Conversely take $m=-k$ and
$n=((\rho,x)+7k)/6$ for a vector in (NL3).
If the image vanishes, $x=m\rho/6$; pairing with $\alpha$
forces $m=6j$, and then $n=7j$, so the kernel is exactly
$\mathbb Zw$. Finally
\[
 |x-m\rho/6|^2
 =x^2-\tfrac m3(\rho,x)+\tfrac73m^2=x^2-2mn.
\]
Polarization gives the full pairing, proving the isometry.

For every $T\in O(N)$ the exact transport laws are
\[
 T(K_\rho)=K_{T\rho},\qquad
 T(\Lambda_\rho)=\Lambda_{T\rho},\qquad
 (x,m,n)\mapsto(Tx,m,n):w_\rho\mapsto w_{T\rho}.
 \tag{NL18}
\]
They follow by substitution in (NL3).
In particular the stabilizer of $\rho$ acts on the fixed
$\Lambda_\rho$. General isometries act on this specified
family of neighbors; no assertion that an arbitrary chosen
automorphism complement fixes this Weyl vector is required.

\subsection{The exact $F_4$ Coxeter action of order twelve}

The order twelve in (NL14) is a lattice quotient invariant of
this explicitly marked construction. There is also the ambient
$F_4$ Coxeter action on the same original $D_4$ real root space.
Use the precise $F_4$ simple roots
\[
 \beta_1=e_2-e_3,\quad\beta_2=e_3-e_4,\quad
 \beta_3=e_4,\quad\beta_4=(e_1-e_2-e_3-e_4)/2,
 \qquad C=s_{\beta_1}s_{\beta_2}s_{\beta_3}s_{\beta_4},
\]
where $s_\beta(x)=x-2(x,\beta)\beta/\beta^2$ and the rightmost
factor acts first. The first two roots have squared norm two,
and the last two squared norm one. No rescaling of the $D_4$
roots has occurred.

\begin{theorem}[Coxeter twelve and its marked triality quotient]
\label{esnl:thm:f4-coxeter-twelve}
In the retained four-coordinate metric,
\[
 C=\frac12
 \begin{pmatrix}
 1&1&1&1\\-1&1&1&-1\\1&1&-1&-1\\1&-1&1&-1
 \end{pmatrix},\qquad
 C^3=
 \begin{pmatrix}
 0&1&0&0\\-1&0&0&0\\0&0&0&-1\\0&0&1&0
 \end{pmatrix}.
 \tag{NL19}
\]
It satisfies $C^4-C^2+I=0$, $C^6=-I$ and has exact order
twelve. Its marked discriminant action is
$v_D+D\mapsto c_D+D\mapsto s_D+D\mapsto v_D+D$.
The restriction of the original triality quotient is the exact
sequence
\[
 1\longrightarrow\langle C^3\rangle\cong\mathbb Z/4
 \longrightarrow\langle C\rangle\cong\mathbb Z/12
 \longrightarrow\langle(v_D\,c_D\,s_D)\rangle\cong\mathbb Z/3
 \longrightarrow1,
 \tag{NL20}
\]
with $\langle C^3\rangle\subset W(D_4)$.
\end{theorem}

\begin{proof}
The first reflection exchanges coordinates two and three; the
second exchanges coordinates three and four; the third changes
the sign of coordinate four; the last is
$I-2\beta_4\beta_4^T$. Their product gives the first matrix in
(NL19). Multiplying it three times gives the second, whose
square is $-I$. Multiplying the first matrix also gives
$C^4-C^2+I=0$, and hence $C^{12}=I$.
Its first column gives
$Cv_D=c_D+(1,-1,0,0)$, with the displayed correction in $D$.
Directly $Cs_D=v_D$ and $Cc_D=(0,0,-1,-1)+s_D$.
The displayed integer correction lies in $D$ and proves the
stated three-cycle with the original $v_D,s_D,c_D$ markings.
The first two reflections preserve $D$ as root reflections.
The last two preserve $D$ because a sign change preserves
integer sum parity, and reflection in $\beta_4$ sends an
integer even-sum vector to an integer vector of even sum.
The inverse reflections have the same property, so these
maps are isometries of $D$.
Consequently their induced action on $D^\vee/D$ is defined
and has the exact order-three image calculated above.
The matrix $C^3$ is an even signed permutation of exact order
four, and thus lies in the already identified group $W(D_4)$.
If $C^m=I$, its discriminant image forces $3\mid m$,
and the exact order of $C^3$ then forces $4\mid m/3$.
Thus $C$ has exact order twelve and the discriminant kernel
on its cyclic subgroup is exactly $\langle C^3\rangle$.
This proves every map and group order in (NL20).
\end{proof}

In particular $C^2$ still has nontrivial order-three
discriminant action, while $C^3$ is in the long-root Weyl
group. The Coxeter denominator six in (NL1), the cyclic order
twelve in (NL14), and the $F_4$ Coxeter twelve in (NL20)
therefore have their stated, distinct domains and explicit
relationships through the original $D_4$ marking.
